\section{Receiver-Owned Bilateral Admission}
\label{sec:design}

The protocol separates evidence produced by the sender from authority retained by the receiver. A sender may supply receipts and a DSSE envelope. It may not supply the receiver's expected keys, treaty record, trust bundle, or continuation state. Those objects are resolved by identifier from receiver-owned stores.

\subsection{Artifacts and bindings}

\paragraph{Receipt.}
A Chio receipt records the capability identifier, tool server and name, canonical action hash, allow or deny decision, policy hash, evidence, metadata, trust level, and kernel public key. The body is signed over canonical JSON. A valid vendor receipt proves only that the corresponding key signed that body. The receiver still has to decide whether the receipt belongs to the current action.

\paragraph{Treaty scope and ladder.}
A \emph{TreatyScope} fixes a treaty identifier, participant kernel identifiers and public keys, hashes of participant ladder manifests, allowed action classes, validity interval, revocation epoch, and trust-bundle digest. Each ladder manifest assigns an enforcement mode and evidence requirement to an action class. The finite order is observation, guarded, receipt-backed, partition-contingency, and maintenance. Intersection takes the stricter participant mode and rejects a destructive class below receipt-backed. An unknown class denies. Each action class also declares a consistency model, one of crdt-commutative, totally-ordered, or quorum-required; quorum-required is a consistency model, not a ladder mode.

\paragraph{Continuation and lineage.}
A \emph{CrossKernelContinuation} links child work to a parent receipt and session anchor, capability, action class, target tool, nonce, source and target kernels, and validity interval. A \emph{ReceiptLineageStatement} binds the parent and child receipt digests to that continuation and to the bilateral invocation. The receiver consumes the continuation during admission evaluation, records the consumed identifier in the accepted result, and releases it only when pre-dispatch execution denies or aborts.

\paragraph{Bilateral predicate.}
The strict predicate type is \codepath{chio.bilateral-cosign-invocation.v1}. Its treaty reference binds the treaty scope and ladder-intersection digests; admission report; continuation; lineage bundle; action class and consistency model; request and outcome digests; local and remote receipt digests; lease and governance references; and participant kernel identifiers. DSSE signs the canonical statement through its pre-authentication encoding (PAE) \cite{dsse}. The envelope contains exactly two signatures whose key identifiers must be distinct and match the configured public keys. Figure~\ref{fig:treaty} shows the artifact.

Two details matter. Both signatures must cover the same predicate bytes; two
signatures over different statements do not compose into bilateral evidence.
Distinct keys also do not establish distinct controllers. The verifier checks
the bytes and key identifiers, while deployment policy is responsible for key
custody.

\begin{figure*}[t]
  \centering
  \resizebox{.83\textwidth}{!}{\begin{tikzpicture}[
  node distance=.8cm and 1.15cm,
  box/.style={draw, rounded corners=2pt, align=center, inner sep=6pt, font=\footnotesize, text width=3.0cm},
  sig/.style={draw, rounded corners=2pt, align=center, inner sep=6pt, font=\footnotesize, fill=black!6, text width=2.7cm},
  store/.style={draw, dashed, rounded corners=2pt, align=center, inner sep=6pt, font=\footnotesize, text width=3.5cm},
  arr/.style={-Latex, line width=.5pt}
]
\node[box] (left) {Participant A\\policy and key};
\node[box, right=of left] (predicate) {Canonical bilateral predicate\\request, outcome, treaty, continuation, lineage, receipts};
\node[box, right=of predicate] (right) {Participant B\\policy and key};
\node[sig, below left=of predicate] (siga) {Signature A\\over DSSE PAE};
\node[sig, below right=of predicate] (sigb) {Signature B\\over identical PAE};
\node[store, below=1.15cm of predicate] (receiver) {Receiver-owned resolution and verification\\expected keys, scope, manifests, leases, governance, receipts};
\node[box, below=of receiver] (verdict) {Pre-dispatch receiver verdict};
\draw[arr] (left) -- (predicate);
\draw[arr] (right) -- (predicate);
\draw[arr] (predicate) -- (siga);
\draw[arr] (predicate) -- (sigb);
\draw[arr] (siga) -- (receiver);
\draw[arr] (sigb) -- (receiver);
\draw[arr] (receiver) -- (verdict);
\end{tikzpicture}
}
  \caption{The bilateral artifact. Participant policy produces a common treaty-bound predicate; two configured keys sign the same DSSE PAE bytes. The receiver then resolves and checks every referenced object from its own trust context.}
  \Description{Two participant policy inputs lead to a common treaty and predicate. Two distinct signatures cover the predicate. A receiver-owned verifier resolves the treaty, receipts, continuation, lineage, lease, and governance evidence before admission.}
  \label{fig:treaty}
\end{figure*}

\subsection{Admission sequence}

Figure~\ref{fig:admission} shows the receiver's control flow.

\begin{enumerate}[noitemsep,topsep=2pt,leftmargin=1.4em]
  \item The kernel first performs its ordinary capability, scope, revocation, and guard checks. Budget admission runs after the treaty hook, so a treaty denial never charges the budget.
  \item Trusted runtime state classifies the request. Local requests follow
  the ordinary local admission path. A request classified as federated without
  complete treaty context denies.
  \item The hook parses only identifiers and bindings from request context. Any attempt to carry a trust root, peer directory, signing key, or dynamic trust bundle in request metadata denies.
  \item Receiver stores resolve the treaty scope, ladder intersection, continuation, lineage bundle, bilateral invocation, bilateral DSSE envelope, and the admission bundle that names the capability lease and governance receipt.
  \item Validators check schema, canonical hashes, freshness, participant sets, action class, consistency model, required evidence, continuation origin and target, and cross-artifact equality.
  \item The hook checks the policy summary, every treaty binding, the lease and governance references against its admission bundle, and the signer set against the treaty participants. The strict envelope verifier then checks the payload type, exactly two signatures, canonical payload bytes, statement and predicate types, one subject, two distinct key identifiers, and both signatures. Peer pins, ladder-manifest freshness, and the revocation epoch are checked by the offline buyer verifier (Section~\ref{sec:implementation}).
  \item Admission either returns a named failure before tool invocation or permits dispatch. On the admitted path, the runtime emits bound receipt metadata and assembles buyer-review evidence.
\end{enumerate}

\begin{figure}[!t]
  \centering
  \resizebox{\columnwidth}{!}{\begin{tikzpicture}[
  node distance=.72cm and .85cm,
  box/.style={draw, rounded corners=2pt, align=center, inner sep=5pt, font=\footnotesize, text width=3.0cm},
  decision/.style={draw, diamond, aspect=2.3, align=center, inner sep=1.5pt, font=\scriptsize, text width=2.1cm},
  allowterm/.style={draw, rounded corners=2pt, align=center, inner sep=5pt, font=\footnotesize, fill=black!7, text width=2.5cm},
  denyterm/.style={draw, double, rounded corners=2pt, align=center, inner sep=5pt, font=\footnotesize, text width=2.5cm},
  arr/.style={-Latex, line width=.5pt}
]
\node[box] (request) {Federated tool request\\after core capability checks};
\node[box, below=of request] (resolve) {Resolve treaty, continuation, lineage, receipts, and keys from receiver stores};
\node[decision, below=of resolve] (bound) {Fresh and mutually bound?};
\node[denyterm, right=of bound, xshift=.65cm] (denybind) {Deny with named binding or freshness code};
\node[box, below=of bound] (verify) {Verify canonical DSSE, two configured keys, policy, lease, and governance};
\node[decision, below=of verify] (admit) {Receiver policy accepts?};
\node[denyterm, right=of admit, xshift=.65cm] (denypolicy) {Deny before dispatch};
\node[allowterm, below=of admit] (dispatch) {Dispatch tool and bind receipt evidence};
\draw[arr] (request) -- (resolve);
\draw[arr] (resolve) -- (bound);
\draw[arr] (bound) -- node[above, font=\scriptsize]{no} (denybind);
\draw[arr] (bound) -- node[left, font=\scriptsize]{yes} (verify);
\draw[arr] (verify) -- (admit);
\draw[arr] (admit) -- node[above, font=\scriptsize]{no} (denypolicy);
\draw[arr] (admit) -- node[left, font=\scriptsize]{yes} (dispatch);
\end{tikzpicture}
}
  \caption{Receiver-side pre-dispatch admission. Trust material flows from receiver-owned stores, not from the request. Every denial terminal precedes the tool boundary.}
  \Description{A federated request passes core authorization, receiver-owned evidence resolution, cross-artifact verification, and receiver policy before tool dispatch. Failures terminate at denial nodes before the tool.}
  \label{fig:admission}
\end{figure}

\paragraph{Receiver-owned policy.}
Explicit Rust validators and the strict bilateral verifier make the hook's
admission decision. The bounded evaluator in Section~\ref{sec:model} operates
on a receipt view constructed after those checks and is not a substitute
authorization path.

\paragraph{Failure evidence.}
The negative matrix assigns each attack a stable identifier, enforcement
phase, and expected diagnostic. Kernel tests and the dispatch-capable
benchmark also inspect the tool invocation counter. Envelope-only cases stop
earlier, before entering an API that can dispatch a tool.
